\chapter{Μεθοδολογια}

\section{Ορισμός του προβλήματος}

Η παρούσα εργασία μελετά την αυτόματη αναγνώριση ανωμαλιών σε δεδομένα καταγραφής του συστήματος \en{HDFS}. Το σύνολο δεδομένων προέρχεται από το \en{Loghub} και αποτελείται από ακολουθίες συμβάντων, οι οποίες αντιστοιχούν σε διαφορετικές καταστάσεις λειτουργίας του συστήματος. Κάθε ακολουθία συσχετίζεται με μία δυαδική ετικέτα, ώστε το πρόβλημα να διατυπώνεται ως ταξινόμηση μεταξύ φυσιολογικής και ανώμαλης συμπεριφοράς.

Έστω

\[
D = \{(x_j, y_j)\}_{j=1}^{M}
\]

το σύνολο εκπαίδευσης, όπου \(x_j\) είναι μία ακολουθία \en{logs} και \(y_j \in \{0,1\}\) η αντίστοιχη ετικέτα. Το ζητούμενο είναι να προσαρμοστεί ένα μοντέλο \(f_{\theta}\), βασισμένο στην αρχιτεκτονική \en{BERT}, ώστε για κάθε είσοδο \(x_j\) να εκτιμά την πιθανότητα εμφάνισης ανωμαλίας.

Η προσαρμογή του μοντέλου πραγματοποιείται με χρήση της μεθόδου \en{LoRA} (\en{Low-Rank Adaptation}), κατά την οποία οι αρχικές παράμετροι του προεκπαιδευμένου μοντέλου παραμένουν σταθερές και εκπαιδεύεται μόνο ένα μικρό σύνολο πρόσθετων παραμέτρων προσαρμογής. Οι εκπαιδεύσιμες αυτές παράμετροι συμβολίζονται συνολικά ως \(\theta\).

Η εκπαίδευση εξετάζεται σε δύο διαφορετικές συνθήκες. Στην πρώτη, τα δεδομένα συγκεντρώνονται σε ένα ενιαίο σημείο και το μοντέλο βελτιστοποιείται πάνω στο πλήρες σύνολο εκπαίδευσης:

\[
\theta^{*}
=
\arg\min_{\theta}
\mathcal{L}(f_{\theta}(x), y)
\]

Στη δεύτερη συνθήκη, το σύνολο δεδομένων χωρίζεται σε επιμέρους τμήματα

\[
D_1,D_2,\ldots,D_N
\]

τα οποία αντιστοιχούν σε διαφορετικούς \en{clients}. Κάθε \en{client} εκπαιδεύει τοπικά το μοντέλο στα δικά του δεδομένα και επιστρέφει ενημερωμένες παραμέτρους στον κεντρικό διακομιστή. Οι παράμετροι που ανταλλάσσονται αντιστοιχούν αποκλειστικά στις εκπαιδεύσιμες παραμέτρους της μεθόδου \en{LoRA} και συμβολίζονται συνολικά ως $\theta$.


Στην περίπτωση της στρατηγικής \en{FedAvg}, ο διακομιστής παράγει το καθολικό μοντέλο μέσω σταθμισμένου μέσου όρου των παραμέτρων που επιστρέφουν οι \en{clients}:

\[
\theta^{t+1}
=
\sum_{k=1}^{N}
\frac{n_k}{\sum_{i=1}^{N}n_i}
\theta_k^{t+1}
\]

όπου $n_k$ είναι το πλήθος δειγμάτων του $k$-οστού \en{client}.

Στην περίπτωση της στρατηγικής \en{FedAdam}, ο διακομιστής υπολογίζει αρχικά τη συνολική ενημέρωση:

\[
g^t
=
\sum_{k=1}^{N}
\frac{n_k}{\sum_{i=1}^{N}n_i}
\left(
\theta_k^{t+1}
-
\theta^t
\right)
\]

και στη συνέχεια εφαρμόζει ενημέρωση τύπου \en{Adam}:

\[
m^t
=
\beta_1m^{t-1}
+
(1-\beta_1)g^t
\]
\[
v^t
=
\beta_2v^{t-1}
+
(1-\beta_2)(g^t)^2
\]
\[
\theta^{t+1}
=
\theta^t
+
\eta
\frac{m^t}
{\sqrt{v^t}+\tau}
\]

όπου $\eta$ είναι ο ρυθμός μάθησης του διακομιστή, ενώ οι παράμετροι $\beta_1$, $\beta_2$ και $\tau$ αντιστοιχούν στους συντελεστές πρώτης και δεύτερης ροπής και στον όρο σταθεροποίησης της μεθόδου αντίστοιχα.

\section{Αρχιτεκτονική του συστήματος}

Η υλοποίηση της παρούσας εργασίας οργανώνεται σε δύο διακριτά αλλά αλληλένδετα έργα λογισμικού. Το πρώτο, \en{hdfs\_anomaly\_detection}, υλοποιεί τη ροή προεπεξεργασίας του συνόλου δεδομένων και την κεντρικοποιημένη εκπαίδευση του μοντέλου. Το δεύτερο, \en{flower\_hdfs\_anomaly\_detection}, επεκτείνει την ίδια βασική ροή σε περιβάλλον ομοσπονδιακής μάθησης, χρησιμοποιώντας το \en{Flower}.

Σε υψηλό επίπεδο, το σύστημα ακολουθεί την εξής ακολουθία επεξεργασίας:
\begin{enumerate}

\item φόρτωση και μετατροπή των δεδομένων \en{HDFS} σε ακολουθίες κειμένου με δυαδικές ετικέτες,

\item διαχωρισμός των δεδομένων σε σύνολα εκπαίδευσης, επικύρωσης και ελέγχου,

\item \en{tokenization} των ακολουθιών με τον \en{tokenizer} του επιλεγμένου μοντέλου,

\item προσαρμογή ενός προεκπαιδευμένου \en{BERT} για δυαδική ταξινόμηση με χρήση \en{LoRA} κεντρικά και σε ομοσπονδιακό περιβάλλον,

\item αξιολόγηση με μετρικές ταξινόμησης και καταγραφή των αποτελεσμάτων.

\end{enumerate}

  

Η ίδια αναπαράσταση δεδομένων χρησιμοποιείται και στις δύο πειραματικές συνθήκες. Στην κεντρικοποιημένη εκπαίδευση το μοντέλο έχει πρόσβαση στο πλήρες σύνολο εκπαίδευσης. Στην ομοσπονδιακή εκπαίδευση, το σύνολο εκπαίδευσης χωρίζεται σε επιμέρους τμήματα, τα οποία ανατίθενται στους \en{clients}, ενώ τα σύνολα επικύρωσης και ελέγχου χρησιμοποιούνται κεντρικά για την αποτίμηση του καθολικού μοντέλου.


\section{Σύνολο δεδομένων \en{Loghub HDFS}}

Για την πειραματική αξιολόγηση χρησιμοποιείται το σύνολο δεδομένων \en{HDFS} της συλλογής \en{Loghub}. Το \en{Loghub} αποτελεί οργανωμένη συλλογή πραγματικών αρχείων καταγραφής από διαφορετικά συστήματα λογισμικού και έχει δημιουργηθεί με σκοπό να υποστηρίξει αναπαραγώγιμη έρευνα στην αυτοματοποιημένη ανάλυση \en{logs} και στην ανίχνευση ανωμαλιών \cite{zhu2023loghub}. Το συγκεκριμένο σύνολο \en{HDFS} προέρχεται από το \en{Hadoop Distributed File System} και βασίζεται στην εργασία των \en{Xu et al.}, οι οποίοι μελέτησαν την ανίχνευση προβλημάτων μεγάλης κλίμακας μέσω εξόρυξης αρχείων καταγραφής κονσόλας \cite{xu2009detecting}.

Το \en{HDFS} είναι ιδιαίτερα κατάλληλο για την παρούσα εργασία, επειδή τα αρχεία καταγραφής του μπορούν να οργανωθούν φυσικά σε ακολουθίες με βάση το αναγνωριστικό του \en{HDFS block}. Σε ένα κατανεμημένο σύστημα αρχείων, η εκτέλεση εργασιών πάνω σε ένα \en{block} παράγει πολλαπλές γραμμές καταγραφής που περιγράφουν διαφορετικά βήματα της ίδιας λειτουργικής ροής. Επομένως, η μονάδα ταξινόμησης δεν είναι η μεμονωμένη γραμμή \en{log}, αλλά ολόκληρη η ακολουθία γεγονότων που αντιστοιχεί σε ένα συγκεκριμένο \en{BlockId}. Στη βιβλιογραφία αυτή η ομαδοποίηση αναφέρεται συχνά ως ομαδοποίηση με βάση το \en{session id}, όπου στο \en{HDFS} το \en{session id} ταυτίζεται πρακτικά με το αναγνωριστικό \en{blk\_} του \en{block}.

Κάθε τέτοια ακολουθία συνοδεύεται από δυαδική ετικέτα. Η ετικέτα \en{Normal} δηλώνει φυσιολογική συμπεριφορά, ενώ η ετικέτα \en{Anomaly} δηλώνει ότι η αντίστοιχη ακολουθία σχετίζεται με προβληματική ή ανώμαλη λειτουργία. Με βάση αυτή τη δομή, το πρόβλημα διατυπώνεται ως επιβλεπόμενη δυαδική ταξινόμηση ακολουθιών καταγραφής.

\subsection{Ανάλυση και ομαδοποίηση των καταγραφών}

\subsection{Δημιουργία του συνόλου εισόδου του μοντέλου}

Μετά την ομαδοποίηση, το ενδιάμεσο σύνολο δεδομένων έχει δύο βασικές στήλες:

\begin{itemize}
    \item \en{text}: την ακολουθία των προτύπων καταγραφής που αντιστοιχεί σε ένα \en{HDFS block},
    \item \en{label}: τη δυαδική ετικέτα της ακολουθίας, με τιμές 0 για φυσιολογική λειτουργία και 1 για ανωμαλία.
\end{itemize}

Το σύνολο αυτό μετατρέπεται σε αντικείμενο \en{Dataset} της βιβλιοθήκης \en{Hugging Face Datasets}. Η στήλη \en{label} δηλώνεται ως κατηγορική μεταβλητή δύο κλάσεων, ώστε να υποστηρίζεται στρωματοποιημένος διαχωρισμός. Η διάσπαση γίνεται σε τρία υποσύνολα: \en{train}, \en{val} και \en{test}. Αρχικά κρατείται το 20\% των δεδομένων εκτός εκπαίδευσης και στη συνέχεια αυτό το τμήμα χωρίζεται ισόποσα σε επικύρωση και έλεγχο. Έτσι, προκύπτει αναλογία 80\% για εκπαίδευση, 10\% για επικύρωση και 10\% για τελική αξιολόγηση. Και τα δύο στάδια διαχωρισμού είναι στρωματοποιημένα ως προς την ετικέτα, ώστε η αναλογία φυσιολογικών και ανώμαλων ακολουθιών να παραμένει σταθερή στα επιμέρους σύνολα.

Στο επόμενο βήμα εφαρμόζεται \en{tokenization} με τον \en{tokenizer} του μοντέλου που χρησιμοποιείται στην εκπαίδευση. Στην παρούσα υλοποίηση το προεπιλεγμένο μοντέλο είναι το \en{bert-base-uncased}. Κάθε κειμενική ακολουθία μετατρέπεται στα πεδία \en{input\_ids} και \en{attention\_mask}. Οι ακολουθίες συμπληρώνονται με \en{padding} ή περικόπτονται με \en{truncation} σε μέγιστο μήκος 512 \en{tokens}, ώστε όλες οι είσοδοι να έχουν συμβατή μορφή για το \en{BERT}.

Το τελικό αποτέλεσμα της διαδικασίας είναι ένα \en{DatasetDict} με τα τρία \en{tokenized splits}. Κάθε δείγμα περιέχει πλέον τα \en{input\_ids}, το \en{attention\_mask} και την ετικέτα \en{label}, μορφοποιημένα ως \en{PyTorch tensors}. Το σύνολο αυτό αποθηκεύεται στον δίσκο και χρησιμοποιείται χωρίς περαιτέρω μετασχηματισμό τόσο από την κεντρικοποιημένη εκπαίδευση όσο και από την ομοσπονδιακή υλοποίηση. Στην ομοσπονδιακή περίπτωση, μόνο το υποσύνολο \en{train} διαμοιράζεται στους \en{clients}, ενώ τα \en{val} και \en{test} παραμένουν κοινά σύνολα αξιολόγησης για το καθολικό μοντέλο.

\section{\en{Fine tuning} με \en{LoRA}}

Η προσαρμογή του προεκπαιδευμένου μοντέλου γίνεται με τη μέθοδο \en{LoRA}, ώστε να μην ενημερώνεται ολόκληρο το \en{BERT}. Αν οι παράμετροι του αρχικού μοντέλου συμβολιστούν με \(\phi_0\), τότε κατά την εκπαίδευση παραμένουν σταθερές και βελτιστοποιούνται μόνο οι πρόσθετες παράμετροι προσαρμογής:

\[
\phi_0^{t+1} = \phi_0^t = \phi_0.
\]

Για ένα γραμμικό στρώμα του \en{Transformer}, έστω ότι ο αρχικός πίνακας βαρών είναι \(W_0 \in \mathbb{R}^{d_{\mathrm{out}} \times d_{\mathrm{in}}}\). Στο πλήρες \en{fine tuning} θα εκπαιδευόταν απευθείας ένας πλήρης πίνακας ενημέρωσης \(\Delta W\). Αντίθετα, το \en{LoRA} εκφράζει αυτή την ενημέρωση ως γινόμενο δύο πινάκων χαμηλού βαθμού:

\[
\Delta W = \frac{\alpha}{r}BA,
\qquad
A \in \mathbb{R}^{r \times d_{\mathrm{in}}},
\qquad
B \in \mathbb{R}^{d_{\mathrm{out}} \times r},
\qquad
r \ll \min(d_{\mathrm{in}}, d_{\mathrm{out}}).
\]

Η παράμετρος \(r\) είναι ο βαθμός της χαμηλόβαθμης προσέγγισης, ενώ το \(\alpha\) ελέγχει την κλίμακα της προσαρμογής. Ο πίνακας που χρησιμοποιείται κατά το εμπρόσθιο πέρασμα γίνεται:

\[
W_{\mathrm{LoRA}}
=
W_0
+
\frac{\alpha}{r}BA.
\]

Για είσοδο \(h \in \mathbb{R}^{d_{\mathrm{in}}}\), η έξοδος του στρώματος είναι:

\[
z
=
W_{\mathrm{LoRA}}h
=
W_0h
+
\frac{\alpha}{r}BAh.
\]

Το βασικό μοντέλο παραμένει αμετάβλητο μέσω του όρου \(W_0h\), ενώ η προσαρμογή εισάγεται μόνο από τον χαμηλόβαθμο όρο \(\frac{\alpha}{r}BAh\). Οι εκπαιδεύσιμες παράμετροι ενός τέτοιου στρώματος είναι μόνο οι πίνακες \(A\) και \(B\), άρα το πλήθος τους είναι:

\[
P_{\mathrm{LoRA}}
=
r(d_{\mathrm{in}} + d_{\mathrm{out}}),
\qquad
\frac{P_{\mathrm{LoRA}}}{P_{\mathrm{full}}}
=
\frac{r(d_{\mathrm{in}} + d_{\mathrm{out}})}
{d_{\mathrm{out}}d_{\mathrm{in}}}.
\]

Στην παρούσα υλοποίηση, η \en{LoRA} προσαρμογή εφαρμόζεται στους πίνακες \en{query} και \en{value} κάθε στρώματος αυτοπροσοχής του \en{BERT}. Για το επίπεδο \(l\) ισχύει:

\[
W_{Q,\mathrm{LoRA}}^{(l)}
=
W_Q^{(l)}
+
\frac{\alpha}{r}B_Q^{(l)}A_Q^{(l)},
\]

\[
W_{V,\mathrm{LoRA}}^{(l)}
=
W_V^{(l)}
+
\frac{\alpha}{r}B_V^{(l)}A_V^{(l)}.
\]

Οι προβολές που χρησιμοποιούνται στον μηχανισμό προσοχής είναι επομένως:

\[
Q^{(l)}
=
H^{(l)}W_{Q,\mathrm{LoRA}}^{(l)},
\qquad
V^{(l)}
=
H^{(l)}W_{V,\mathrm{LoRA}}^{(l)}.
\]

Ο μηχανισμός αυτοπροσοχής παραμένει ο ίδιος:

\[
\operatorname{Attention}(Q,K,V)
=
\operatorname{softmax}
\left(
\frac{QK^{T}}{\sqrt{d_k}}
\right)
V,
\]

όπου \(d_k\) είναι η διάσταση των διανυσμάτων \en{key}. Επομένως, το \en{LoRA} δεν αλλάζει τον τύπο της προσοχής, αλλά αλλάζει αποδοτικά τις προβολές \(Q\) και \(V\).

\section{Κεντρικοποιημένη Εκπαίδευση}

Η κεντρικοποιημένη αρχιτεκτονική υλοποιείται στο \en{hdfs\_anomaly\_detection}. Το σενάριο εκπαίδευσης φορτώνει το αποθηκευμένο \en{tokenized dataset}, χρησιμοποιεί το σύνολο \en{train} για βελτιστοποίηση και το σύνολο \en{val} για αξιολόγηση ανά εποχή. Η εκπαίδευση βασίζεται στον \en{Trainer} της βιβλιοθήκης \en{Transformers}, με προσαρμοσμένη κλάση \en{CustomTrainer} για τον υπολογισμό της συνάρτησης απώλειας.


Η προσαρμοσμένη συνάρτηση απώλειας υποστηρίζει δύο μηχανισμούς για την αντιμετώπιση της ανισορροπίας των κλάσεων. Ο πρώτος είναι η χρήση σταθμισμένης \en{cross entropy}, όπου τα βάρη υπολογίζονται από την κατανομή των ετικετών στο σύνολο εκπαίδευσης. Ο δεύτερος είναι η χρήση \en{focal loss}, η οποία δίνει μεγαλύτερη έμφαση σε δυσκολότερα παραδείγματα. Οι επιλογές αυτές ενεργοποιούνται παραμετρικά από το αρχείο ρυθμίσεων.
  
Το αποτέλεσμα της κεντρικοποιημένης εκπαίδευσης είναι ένα σύνολο καταγραφών που περιλαμβάνει τις υπερπαραμέτρους του πειράματος και την ιστορική εξέλιξη των μετρικών. Το υποσύστημα αυτό χρησιμοποιείται για τη διερεύνηση βασικών επιλογών, όπως ο ρυθμός μάθησης, το μέγεθος \en{batch}, η τάξη του \en{LoRA} και οι επιλογές απώλειας, πριν από την εκτέλεση των ομοσπονδιακών πειραμάτων.

\section{Ομοσπονδιακή Εκπαίδευση}
Η ομοσπονδιακή αρχιτεκτονική υλοποιείται στο \en{flower\_hdfs\_anomaly\_detection} και βασίζεται στο \en{Flower}. Αποτελείται από δύο κύρια μέρη, το \en{ServerApp} και \en{ClientApp}. Ο \en{server} έχει ρόλο συντονιστή: αρχικοποιεί το μοντέλο, εξάγει τις παραμέτρους \en{LoRA}, επιλέγει τη στρατηγική συνάθροισης και οργανώνει τους γύρους εκπαίδευσης. Οι \en{clients} εκτελούν τοπική εκπαίδευση πάνω στα δικά τους τμήματα του συνόλου δεδομένων και επιστρέφουν μόνο τις ενημερωμένες παραμέτρους του μοντέλου μαζί με βασικές μετρικές εκπαίδευσης.

  

Η διανομή των δεδομένων στους \en{clients} γίνεται μέσω \en{Dirichlet partitioning} πάνω στο σύνολο εκπαίδευσης και με βάση την στήλη \en{label}. Η παράμετρος $\alpha$ ελέγχει τον βαθμό ανομοιογένειας μεταξύ των τοπικών συνόλων. Με μεγαλύτερες τιμές η κατανομή των κλάσεων στους \en{clients} πλησιάζει περισσότερο την συνολική κατανομή του συνόλου δεδομένων, ενώ με μικρότερες τιμές δημιουργούνται περισσότερο μη ομοιόμορφα τοπικά σύνολα. Έτσι, η αρχιτεκτονική μπορεί να προσομοιώσει συνθήκες \en{non-IID} δεδομένων, οι οποίες είναι χαρακτηριστικές σε πραγματικά κατανεμημένα περιβάλλοντα.

  

Σε κάθε γύρο, ο \en{server} αποστέλλει την τρέχουσα καθολική κατάσταση των \en{LoRA} παραμέτρων στους επιλεγμένους \en{clients}. Κάθε \en{client} δημιουργεί τοπικά το ίδιο μοντέλο \en{BERT} με \en{LoRA}, φορτώνει τις παραμέτρους που έλαβε και εκτελεί εκπαίδευση με τον \en{CustomTrainer}. Μετά το τέλος της τοπικής εκπαίδευσης, ο \en{client} επιστρέφει την ενημερωμένη κατάσταση \en{PEFT} και το πλήθος των παραδειγμάτων που χρησιμοποιήθηκαν. Ο \en{server} συνδυάζει τις ενημερώσεις με στρατηγικές όπως \en{FedAvg} ή \en{FedAdam}, παράγοντας το νέο καθολικό μοντέλο για τον επόμενο γύρο.

  

Η αξιολόγηση του καθολικού μοντέλου πραγματοποιείται κεντρικά μετά από κάθε γύρο, χρησιμοποιώντας το σύνολο επικύρωσης. Το σύστημα καταγράφει το \en{validation loss}, το \en{accuracy}, το \en{precision}, το \en{recall} και το \en{F1-score}. Στο τέλος της εκπαίδευσης μπορεί να εκτελεστεί επιπλέον \en{threshold tuning}. Σε αυτή τη διαδικασία, το μοντέλο παράγει πιθανότητες για την κλάση της ανωμαλίας, αναζητείται το \en{threshold} που μεγιστοποιεί το \en{F1-score} στο σύνολο επικύρωσης και το επιλεγμένο \en{threshold} εφαρμόζεται στο σύνολο ελέγχου.

\section{Μετρικές Αξιολόγησης}
